\chapter{Prefácio}

Este livro tem uma ambição simples: uma única obra coerente de ciências
da Terra e do espaço, que cubra tudo do primeiro ano da escola ao fim da
graduação, para leitores de qualquer lugar do mundo.

O estilo é conciso e rigoroso: cursos construídos a partir de
definições, exemplos, proposições, métodos e esquemas, com um raciocínio
cuidadoso sempre que for acessível ao nível considerado, seguidos de
exercícios cuidadosamente escolhidos. Os resultados enunciados sem
dedução completa são marcados explicitamente como admitidos. As soluções
completas de todos os exercícios estão reunidas no final do livro. Cada
número citado sobre a Terra real, os seus oceanos, a sua atmosfera ou o
céu vem de uma fonte de referência, e cada curva desenhada a partir de
medidas reais é calculada a partir dos dados publicados.

O curso se divide em cinco livros. O primeiro cobre os anos do ensino
fundamental, o segundo os três anos do ensino médio, e os três últimos os
três anos da universidade. Os dois livros escolares nunca ensinam a mesma
noção duas vezes: cada ideia é ensinada uma só vez, no ano em que ela
cabe pela primeira vez, e um capítulo posterior que precise dela começa
com um breve lembrete, \emph{O que você já sabe}, antes de ir adiante.
Quando uma noção pertence à física, à biologia ou à química, os livros
irmãos da mesma coleção a ensinam, e este livro a relembra antes de
adotar o ponto de vista de quem estuda a Terra. Nos livros escolares, cada
ano é uma parte, dividida em capítulos; cada livro universitário é um
ano, dividido em capítulos. Quanto mais cedo o ano, mais jovem o leitor
para quem ele é escrito: mais ilustrações, etapas resolvidas mais
detalhadas e exercícios mais suaves.

\section*{Como ler este livro}

Os enunciados são numerados dentro de cada capítulo e identificados por
cores: as definições em azul, os teoremas em vermelho, as proposições e os
seus parentes em laranja, e os métodos --- receitas curtas para tarefas
comuns --- em verde. Os exercícios são graduados de ($\star$), para as
aplicações diretas, a ($\star\star\star$), para os problemas mais
difíceis. Os quadros emoldurados trazem o que não faz parte do raciocínio:
um lembrete do que um capítulo anterior ensinou, como os geólogos e os
outros cientistas da Terra trabalham no campo ou no laboratório, uma
página de história, dados que mudam de ano para ano (com o ano a que se
referem) e conselhos de segurança.

\section*{Segurança}

A Terra se aprende olhando para ela: o céu, o tempo que faz, uma
paisagem, as pedras de um prédio, o solo de um jardim. É tudo o que este
livro pede ao seu leitor. Nunca olhe para o Sol, com ou sem instrumento:
ele cega. Pedreiras, falésias, cavernas, rios em cheia e vulcões são
lugares para profissionais com formação e equipamento, e este livro
descreve o trabalho deles sem mandar ninguém para lá.

\section*{Contribuir}

O código-fonte deste livro fica num repositório git público:
\begin{center}
\url{https://github.com/bvirrion/one-earth-and-space-sciences-book}
\end{center}
As contribuições --- novos capítulos, correções, explicações melhores,
exercícios adicionais --- são bem-vindas; veja o arquivo
\texttt{CONTRIBUTING.md} na raiz do repositório.
